\chapter{SOURCE CODE}
\label{app:code}

This appendix lists the three places where a design decision is embedded in the
code rather than in prose. The complete repository, including the training,
evaluation, partitioning and reporting scripts together with the Colab notebook
used to run them, is available at
\url{https://github.com/ErlindSkura/Erlind_Skura_Thesis}.

The listings are trimmed to their substantive parts, since imports and argument
parsing add pages without adding information.

\section*{A.1 \quad Polygon Area and Physical Calibration}

The object area is computed analytically from the polygon vertices rather than
from a rasterised mask, for the reason given in Section~\ref{sec:parsing}.

\begin{lstlisting}[language=Python, caption={Calibration constants and shoelace area (\texttt{code/config.py}, \texttt{code/prepare\_data.py})}]
# Calibrated from the burned-in scale bars: the field width is 284.4 um
# at 1000x and scales inversely with magnification, over 1024 px.
NM_PER_PX = {500: 555.6, 1000: 277.8, 3000: 92.6}

BANNER_H = 32  # burned-in Zeiss info bar at the bottom of every micrograph


def polygon_area(points):
    """Shoelace area of a closed polygon, in px^2."""
    n = len(points)
    return abs(
        sum(
            points[i][0] * points[(i + 1) % n][1]
            - points[(i + 1) % n][0] * points[i][1]
            for i in range(n)
        )
    ) / 2.0
\end{lstlisting}

\section*{A.2 \quad Augmentation Geometry}

Augmentation transforms the polygon vertices and the image together and
rasterises only afterwards, so overlapping beads remain separable. Vertices
are continuous coordinates, so a reflection is $x \mapsto W - x$ and not
$x \mapsto W - 1 - x$. The check below is the one that found that error
(Section~\ref{sec:challenges}): it rasterises the transformed polygons and
compares them against the directly transformed raster mask, and the two agree to
an intersection over union of 0.9995 under the correct convention and of only
0.87 under the incorrect one.

\begin{lstlisting}[language=Python, caption={Vertex-space reflection and the check that verifies it (\texttt{code/data\_io.py})}]
def flip_h(img, polys):
    w = img.size[0]
    return (img.transpose(Image.FLIP_LEFT_RIGHT),
            [np.column_stack([w - p[:, 0], p[:, 1]]) for p in polys])


for name, transform in [("flip_h", flip_h), ("flip_v", flip_v),
                        ("rot90", rot90)]:
    moved_img, moved_polys = transform(image, record.polys)
    from_vertices = rasterise(moved_polys, *moved_img.size).any(0)
    from_raster = np.asarray(
        transform(mask_as_image, [])[0].convert("L")) > 127
    print(name, iou(from_vertices, from_raster))
\end{lstlisting}

\section*{A.3 \quad Instance-Level Agreement}

The listing below gives the Aggregated Jaccard Index together with the merge and
split rates that make the counting failure modes explicit. Intersections are
computed only for pairs whose bounding boxes overlap, since a micrograph can hold
152 annotated beads and several hundred predictions, so evaluating every
pair over the full frame would be prohibitive while almost every pair is
disjoint.

\begin{lstlisting}[language=Python, caption={AJI and merge/split rates (\texttt{code/metrics.py})}]
def aji(gts, preds):
    """Aggregated Jaccard Index (Kumar et al., 2017)."""
    inter, iou, ga, pa = iou_matrix(gts, preds)
    union = ga[:, None] + pa[None, :] - inter
    c = u = 0.0
    used = set()
    for i in range(len(gts)):
        j = int(iou[i].argmax())
        if iou[i, j] > 0:
            c += inter[i, j]
            u += union[i, j]
            used.add(j)
        else:
            u += ga[i]
    u += sum(pa[j] for j in range(len(preds)) if j not in used)
    return float(c / u)


def merge_split(gts, preds, thr=0.5):
    """A merge is one prediction holding half or more of two or more
    ground-truth beads; a split is one ground-truth bead holding half or
    more of two or more predictions."""
    inter, _, ga, pa = iou_matrix(gts, preds)
    cov_gt = inter / np.maximum(ga[:, None], 1.0)
    cov_pr = inter / np.maximum(pa[None, :], 1.0)

    merged = set()
    for j in range(len(preds)):
        absorbed = np.nonzero(cov_gt[:, j] >= thr)[0]
        if len(absorbed) >= 2:
            merged.update(absorbed.tolist())

    n_split = sum(1 for i in range(len(gts))
                  if np.count_nonzero(cov_pr[i] >= thr) >= 2)
    return {"merge_rate": 100.0 * len(merged) / len(gts),
            "split_rate": 100.0 * n_split / len(gts)}
\end{lstlisting}
